\chapter{ПРОГРАМНЕ ЗАБЕЗПЕЧЕННЯ}

%Hidden cite for correct bibliography ordering
\nocite{bahvalov-et-al,benerdge-et-al} 

\section{Огляд проблеми}

Для отримання, валідації і оформлення наукових результатів було створене середовище із наступними властивостями:
\begin{itemize}
    \item Легке встановлення на кількох різних машинах.
    \item Контрольовані зміни до налаштувань, які легко поширити на всі машини.
    \item Зручна візуалізація даних.
    \item Можливість працювати із абстракціями близькими до математичного апарату.
\end{itemize}

\noindent Стабільне середовище, яке дозволяє проводити дослідження як на потужних так і на мобільних машинах, дозволяє 
використати час більш продуктивно. З іншої сторони, наявність інтерфейсу високого рівня, який дозволяє працювати із 
формулами майже у співвідношенні один до одного зменшує час на пошуки помилок і відлагодження програмного забезпечення. 

\section{Середовище}

Технологія Docker дозволяє описувати середовище виконання у вигляді коду. Це середовище, яке називається Docker 
відображенням, можна завантажити на спеціальний сервер. На довільній машині із операційною системою Ubuntu можна 
вивантажити це відображення і запустити відповідний контейнер. Він буде містити ідентичне середовище, яке буде 
виконуватися ідентично (при умові, що версії ядер Ubuntu однакові). Також Docker відображенням можна збудувати на 
довільній машині використовуючи текстовий файл опису Dockefile. Це суттєва перевага цієї технології над іншими. Завдяки
таким властивостям можна завантажити Dockerfile у систему зберігання коду і відслідковувати зміни. Це дозволяє будувати
Docker відображення на локальній машині із Dockefile, у випадку, коли немає сервера із потрібним Docker відображенням.

На (рис. \ref{fig:software_docker_hierarchy}) зображена побудована ієрархія Docker відображень. Вона доволі гнучка, бо
дозволяє використовувати відображення із потрібним набором програмного забезпечення на конкретній машині.

\begin{figure}[ht!]
    \centering
    \begin{tikzpicture}
        \begin{class}[text width=2cm]{BaseUI}{0,0}
        \end{class}

        \begin{class}[text width=2cm]{NGSolve}{0,-2}
            \inherit{BaseUI}
        \end{class}

        \begin{class}[text width=2cm]{NGBem}{0,-4}
            \inherit{NGSolve}
        \end{class}

        \begin{class}[text width=2cm]{LaTeX}{0,-6}
            \inherit{NGBem}
        \end{class}

    \end{tikzpicture}
    \caption{Ієрархія Docker відображень}
    \label{fig:software_docker_hierarchy}
\end{figure}


\section{Програмні продукти}

Для виконання обчислень використовувалася програмний пакет NGSolve \cite{software-ngsolve} і його розширення NGBem.

\section{Висновок}

Засоби дозволяють все зберігати у текстових файлах.
